\documentclass[10pt,a4paper]{article}
\usepackage{amsmath,amssymb,amsthm,graphicx}
\usepackage{xfrac}
\newenvironment{enumeratei}
{\bgroup\def\theenumi{\roman{enumi}}\def\theenumii{\arabic{enumii}}\begin{enumerate}}
{\end{enumerate}\egroup}

\hyphenation{addi-tion-al belong belongs denote exchan-ge exchan-ged exchan-ges}

\newcommand{\psfrac}[2]{\sfrac{(#1)}{#2}}

\newcommand\mto{\mathchoice{\longmapsto}{\mapsto}{\mapsto}{\mapsto}}
\let\oldbigsqcup\bigsqcup
\renewcommand{\bigsqcup}{\mathop{\textstyle\oldbigsqcup}\displaylimits}

\let\epsilon\varepsilon

\DeclareMathOperator{\Ess}{Ess}
\DeclareMathOperator{\Leb}{Leb}

\usepackage{tikz}
\usetikzlibrary{plotmarks}

\theoremstyle{plain}
\newtheorem{theorem}{Theorem}[section]
\newtheorem{lemma}[theorem]{Lemma}
\newtheorem{proposition}[theorem]{Proposition}
\newtheorem{corollary}[theorem]{Corollary}
\newtheorem{claim}{Claim}

\theoremstyle{definition}
\newtheorem{remark}[theorem]{Remark}
\newtheorem{example}[theorem]{Example}

\newcommand{\A}{\mathcal{A}}
\newcommand{\B}{\mathcal{B}}
\newcommand{\I}{\mathcal{I}}

\newcommand{\Z}{\mathbb{Z}}
\newcommand{\N}{\mathbb{N}}
\newcommand{\R}{\mathbb{R}}
\title{An interval exchange transformation all of whose connection times are finite is periodic}
\author{Przemyslaw Berk, Frank Trujillo}
\date{doi:10.5802/jep.302; CC BY4.0}
\begin{document}
\maketitle
\noindent\textit{Editorial note.} The counted statement is the author's labelled lemma environment at source0.tex lines 191-193; its proof environment opens at line 195, the line after the blank separator, so the pairing is by adjacency and the retained proof range 195-211 is the whole environment.\par\medskip
\noindent\textit{Editorial note.} This source is new to the material-authoring work: no card had been authored from it before this batch, and its preamble is built from the author's own macro and theorem block (source lines 8-45) rather than from a prepared adapter. The publisher document class, its date commands and the front matter are excluded.\par\medskip
\noindent\textit{Editorial note.} The retained preamble is loaded with the xfrac package, which the author's macro block uses through the sfrac command.\par\medskip
\noindent\textit{Editorial note.} No retained line refers to a label, citation or figure outside the retained span, so the package declares no external reference.\par\medskip
\noindent\textit{Editorial scope.} One original result of the pinned publisher article, counted once: finiteness of the connection times forces periodic behaviour of the interval exchange transformation (author label lem: notallconnections).\par\medskip
\noindent\textit{Allowed standard prerequisites (editorial).}\begin{itemize}
\item Interval exchange transformation. An IET is a piecewise translation of an interval determined by a permutation of the subintervals; the endpoints of the subintervals are its discontinuities.. It is the map studied in the statement.
\item Connection time. M(beta) is the first return time of the left endpoint of the subinterval beta to the set of discontinuities.. The hypothesis is that all these times are finite, which is what makes the orbit set of discontinuities finite.
\item Poincare recurrence. For a finite measure preserving transformation every measurable set of positive measure recurs.. It produces the return time N of the partition element.
\end{itemize}\medskip
\begin{lemma}\label{lem: notallconnections}
Let $T:I\to I$ be an interval exchange transformation given by a permutation $\pi$ and length vector $\lambda$. If for every $\beta\in \A\setminus \{\pi_0^{-1}(1)\}$ we have $M(\beta)<\infty$, then $T$ has only periodic orbits. {More precisely, the base interval $I$ can be decomposed in a finite number of periodic components, given by semi-closed intervals, such that the period is uniform on each of these components.}
\end{lemma}

\begin{proof}
Note that by assumption the set of points
\[
\{T^{n}(\partial I_{\alpha})\mid n\in\Z,\quad\alpha\in \A\}=\{T^{-n}(\partial I_{\alpha})\mid \alpha\in\A\text{ and } 0\le n\le M(\alpha)\}
\]
is finite. Consider the partition given by those points and let $[a,b)$ be an element of this partition.

Note that $T^n$ acts continuously on $[a,b)$ for all $n\in\N$. Indeed, the only possible points of discontinuity of $T$ are $\{\partial I_{\alpha}\}_{\alpha\in\A}$, hence, if for some $n\in\N$ the map $T^n$ did not act continuously on $[a,b)$, there would exist $\beta\in\mathcal A$ such that $\partial I_{\beta}\in T^{n-1}\big([a,b)\big)$, which contradicts the choice of $[a,b)$. In particular, it~follows that $T^n\big([a,b)\big)$ is an interval, for any $n \in \N$.

By Poincaré's recurrence theorem, there exists $N\in\N$ such that
\[
T^{N}\big([a,b)\big)\cap [a,b)\neq \emptyset.
\]
This implies that $T^{N}\big([a,b)\big)=[a,b)$. Indeed, otherwise either $T^{-N}(a)$ or $T^{N}(a)$ belongs to $[a,b)$. Since $a$ is in the orbits of one of the points $\{\partial I_{\alpha}\}_{\alpha\in\A}$, this yields a contradiction.

To sum up, $T^{N}\big([a,b)\big)=[a,b)$ and $T^N$ acts continuously on $[a,b)$. Since $T$ is a piecewise translation, then so is $T^N$. Thus $T^N|_{[a,b)}$ is the identity map on $[a,b)$, which finishes the proof.
\end{proof}

\end{document}
